\section{Investigación y análisis de alternativas}
\label{sec:investigacion}

Antes de escribir código se investigaron cinco problemas con fuentes
contrastables. Esta sección resume las alternativas evaluadas y la
recomendación adoptada. El desarrollo completo está en el repositorio, en
\codigo{docs/research.md}.

\subsection{Representación del importe}

\begin{table}[htbp]
\centering
\caption{Alternativas para representar el dinero.}
\label{tab:dinero}
\small
\begin{tabularx}{\textwidth}{@{}lXX@{}}
\toprule
\textbf{Opción} & \textbf{Ventajas} & \textbf{Limitaciones} \\
\midrule
\codigo{BIGINT} en céntimos & Aritmética exacta y rápida; la invariante
  \enquote{suma de asientos igual a cero} es exacta por construcción &
  No admite fracciones de céntimo; obliga a convertir en la frontera \\
\codigo{NUMERIC(18,2)} & Legible en SQL, escala explícita, exacto &
  Aritmética más lenta; riesgo de deriva de escala \\
\codigo{MONEY} & Tipo dedicado & Precisión dependiente de \codigo{lc\_monetary}
  y de la configuración regional \\
\codigo{FLOAT} & Ninguna aquí & Descartado por la documentación del
  motor~\cite{pg16money} y por~\cite{goldberg1991} \\
\bottomrule
\end{tabularx}
\end{table}

\textbf{Recomendación adoptada:} enteros de céntimos en \codigo{BIGINT}.
\codigo{NUMERIC} es una opción igualmente defendible; se descarta por preferir
que el cuadre global sea cierto en aritmética entera.

\subsection{Fuente de verdad del saldo}

Se evaluaron tres opciones: una columna mutable sin histórico, un libro mayor
puro con el saldo calculado al vuelo, y la combinación de ambos.

La columna mutable sola no deja rastro de auditoría: un error sobrescribe el
saldo sin dejar huella y el cuadre no es verificable de forma independiente.
El libro mayor puro es auditable, pero la lectura del saldo es $O(n)$
creciente y aplicar la no negatividad de forma atómica exige agregar el
histórico bajo bloqueo en cada transferencia.

\textbf{Recomendación adoptada:} libro mayor inmutable como fuente de verdad
más saldo materializado, escritos en la misma transacción, con una prueba de
reconciliación que recalcula desde el libro mayor. Es el patrón que describe
Fowler~\cite{fowler2006} y el que Square adoptó tras observar
inconsistencias en producción~\cite{square2019}.

\subsection{Control de concurrencia}

\begin{table}[htbp]
\centering
\caption{Estrategias de control de concurrencia evaluadas.}
\label{tab:concurrencia}
\small
\begin{tabularx}{\textwidth}{@{}lXX@{}}
\toprule
\textbf{Estrategia} & \textbf{Ventajas} & \textbf{Limitaciones} \\
\midrule
Pesimista con orden determinista &
  Predecible, sin reintentos, deadlocks eliminados por el orden &
  Serializa las operaciones sobre una cuenta caliente \\
\codigo{SERIALIZABLE} con reintento &
  Más concurrencia si el conflicto es raro; garantía formal &
  Exige lógica de reintento; bajo contención alta la tasa de abortos
  crece~\cite{ports2012} \\
\codigo{UPDATE} condicional atómico &
  Muy simple; el motor garantiza la atomicidad de la fila &
  Mezcla validación con escritura; encaja mal con el libro mayor \\
Cola serializada por cuenta &
  Elimina la contención por diseño &
  Exige broker o tabla de cola; complejidad no justificada \\
Cerrojo de Python &
  Ninguna aquí &
  No funciona entre procesos ni entre réplicas. Descartada \\
\bottomrule
\end{tabularx}
\end{table}

\textbf{Recomendación adoptada:} implementar las dos primeras tras una opción
de configuración y \emph{medirlas}. La decisión se detalla en la
sección~\ref{sec:transacciones} y los resultados en
la~\ref{sec:rendimiento}.

\subsection{Arquitectura}

Se compararon monolito tradicional, monolito modular, microservicios,
arquitectura orientada a eventos y hexagonal. El argumento decisivo no es de
moda sino de garantías: \textbf{la transferencia debe ser atómica}. En un
monolito con una base de datos, la atomicidad es un \codigo{COMMIT}. Si se
separan cuentas y libro mayor en servicios con bases distintas, la atomicidad
desaparece y hay que reintroducirla con Saga o un commit en dos fases, lo que
supone compensaciones, estados intermedios visibles y mucha más superficie de
error. Partir el sistema aquí no es una mejora: es cambiar una garantía fuerte
y gratuita por una débil y cara.

Fowler sostiene además que no se debe empezar un proyecto nuevo con
microservicios aunque se prevea que será grande, porque establecer fronteras
correctas es difícil sin patrones de uso reales~\cite{fowler2015}.

\textbf{Recomendación adoptada:} monolito modular con el núcleo de dominio
aislado.

\subsection{Inteligencia artificial}

Se evaluaron reglas deterministas de riesgo, detección de anomalías no
supervisada, clasificación supervisada de fraude, modelos fundacionales y
recuperación aumentada con base de datos vectorial.

Un clasificador supervisado necesitaría transacciones reales etiquetadas, que
no existen y no se pueden fabricar sin falsear el resultado. Sculley
\emph{et al.} documentan que el código de modelo es una fracción pequeña de un
sistema de aprendizaje automático y que el resto es deuda de mantenimiento
continua~\cite{sculley2015}; incurrir en ella sin datos ni evaluación creíble
no aporta nada demostrable.

\textbf{Recomendación adoptada:} ninguna IA en el MVP, y un límite permanente
de diseño: ningún modelo generativo ni estadístico aprueba, rechaza ni
modifica una transferencia o un saldo. Un sistema no determinista no puede
estar en la ruta crítica del dinero.
